\input{../../../../templates/es/preambulo.tex}

\setbeamertemplate{navigation symbols}{}

\setbeamertemplate{footline}[frame number]

\date{}

\begin{document}

\tituloleccion{03}{Selección}

\begin{frame}{Población y pregunta}

Diez candidatos tienen aptitudes ordenadas de $-5$ a $4$. La selección devuelve diez copias de estos mismos candidatos. ¿Quién se reproduce, quién desaparece y cuántas copias conserva el mejor?

\end{frame}

\begin{frame}{Probabilidad proporcional}

Con pesos no negativos $w_i$, $p_i=w_i/\sum_j w_j$ y las copias esperadas son $Np_i$. Como las aptitudes incluyen valores negativos, se añade un piso. Ese piso cambia $p_i$ y por tanto la presión.

\end{frame}

\begin{frame}{Rango y escala}

La selección por rango asigna pesos por posición, no por la distancia entre aptitudes. Una transformación creciente de aptitud conserva el orden y las probabilidades por rango. Aun así, una muestra puede perder al mejor.

\end{frame}

\begin{frame}{Elitismo y SUS}

Elitismo copia intacto al campeón antes de llenar los demás lugares; con aptitud fija conserva al mejor. El muestreo universal estocástico (SUS) usa puntos equiespaciados: misma esperanza que ruleta y menor variación de copias.

\end{frame}

\begin{frame}{Torneo como perilla}

En un torneo con reemplazo de tamaño $k$, elijo el mejor entre $k$ sorteos. Al crecer $k$, aumentan las copias esperadas del campeón y suele caer la diversidad. $k=3$ es una configuración inicial, no una respuesta universal.

\end{frame}

\begin{frame}{Medición antes de elegir}

Registro copias del mejor, individuos extintos, aptitud media y dispersión de genes. En la secuencia extensa de 13 configuraciones, presión y aptitud media correlacionan $+0.95$; presión y dispersión, $-0.92$. Son observaciones del caso, no leyes.

\end{frame}

\begin{frame}{Aplicación a Planta Física}

La selección no repara planes inviables: el evaluador o una reparación documentada deben hacerlo. Para comparar $k$, mantengo el mismo evaluador de 48 decisiones, las mismas semillas y el mismo presupuesto de evaluaciones.

\end{frame}

\begin{frame}{Salida de clase}

Decido una selección inicial y justifico su presión. Registro un riesgo de pérdida de diversidad y la métrica con la que lo vigilaré. Conservo el mejor plan observado fuera de la población si no garantizo elitismo.

\end{frame}

\end{document}
